\documentclass[11pt]{letter}
\usepackage[margin=1in]{geometry}
\usepackage{hyperref}
\usepackage{xcolor}

\signature{Zach DeBruine, Ph.D.\\
Department of Information Sciences and Technology\\
Grand Valley State University\\
Allendale, MI\\
\texttt{zach@singlet.bio}}

\address{Zach DeBruine, Ph.D.\\
Department of Information Sciences and Technology\\
Grand Valley State University\\
Allendale, MI 49401}

\begin{document}

\begin{letter}{Editors\\
\textit{Genome Biology}\\
BMC, Springer Nature}

\opening{Dear Editors,}

We are pleased to submit our manuscript, ``SinglePress: A Purpose-Built File Format for Single-Cell Omics Matrices,'' for consideration as a Software article in \textit{Genome Biology}.

Single-cell RNA-seq datasets are among the most compressible structures in genomics, yet the field stores them in general-purpose formats that waste 80--90\% of disk space.
SinglePress exploits the extreme value concentration of count data---where a median of 72\% of nonzeros equal 1---achieving a median 9.5$\times$ compression over int32 CSC at 868~MB/s decode throughput.
Across 3,253 scRNA-seq datasets spanning nine species and nine protocols, \texttt{.1pz} files are 2--4$\times$ smaller than H5AD and 2.5$\times$ smaller than BPCells, while reading 3--6$\times$ faster.

Three aspects of this work may be of particular interest to the \textit{Genome Biology} readership:

\begin{enumerate}
\item \textbf{Near-optimal compression.} A compression frontier analysis demonstrates that SinglePress approaches or beats the Shannon entropy lower bound.
File size is predictable from nonzero count alone ($R^2 = 0.990$, 0.807 bytes/nonzero), enabling storage cost modeling without trial compression.

\item \textbf{The compression--speed paradox resolved.} Counter to conventional wisdom, better compression yields \emph{faster} reads.
An ablation study decomposes this into a 3.4$\times$ zstd-vs-gzip codec advantage and a 2.1$\times$ threading speedup from the chunked architecture---independent factors that multiply to 3--6$\times$ end-to-end.

\item \textbf{GPU training throughput.} Native PyTorch dataloaders shift deep learning workloads from I/O-bound to compute-bound, delivering 2.5--3$\times$ higher GPU utilization than H5AD-backed alternatives.
As single-cell foundation models scale to billions of cells, I/O efficiency becomes a prerequisite rather than an optimization.
\end{enumerate}

SinglePress provides Python, R, and PyTorch APIs with native interoperability with AnnData, Seurat, SingleCellExperiment, and standard sparse matrix libraries.
The software is open source (GPL-3.0) at \url{https://github.com/Singlet-AI/singlepress}, with format specification, documentation, and all benchmark code included.

This manuscript has not been submitted elsewhere and all authors have approved the submission.
We declare that ZD is a co-founder of Singlet Bio, Inc., which may distribute software incorporating the SinglePress format.

We suggest the following potential reviewers who have expertise in single-cell data infrastructure:
\begin{itemize}
\item Ben Parks (Stanford University) --- developer of BPCells
\item Isaac Virshup (Helmholtz Munich) --- core developer of AnnData/Scanpy
\item Rob Patro (University of Maryland) --- developer of alevin-fry/simpleaf
\end{itemize}

Thank you for your consideration.

\closing{Sincerely,}

\end{letter}
\end{document}
